\section{Proof of Theorem~\ref{thm:main}}\label{sec:proof}

\subsection{The criterion}
We prove Proposition~\ref{prop:criterion}. The proof shows where each of the three properties of $Q_N$
enters. The nonvanishing and the smallness are used once, at the end. The divisibility by squares of primes
is what makes the choice of the prime free.

\begin{proof}[Proof of Proposition~\ref{prop:criterion}]
Let $a/b$ be a rational number with $b$ large. Put $L=\log b$, $U=\kappa L$ and $V=(\kappa-\omega)L$, and
suppose $|\xi-a/b|\le b^{-\rho\kappa}$. By the prime number theorem, $(1-\delta)t\le\theta(t)\le(1+\delta)t$
for all large $t$, and $V\to\infty$ with $b$. If every prime $p\in(V,U]$ had $p^2\mid b$, then
$2(\theta(U)-\theta(V))\le\log b=L$. But
\[
 2\bigl(\theta(U)-\theta(V)\bigr)\ge2L\bigl((1-\delta)\kappa-(1+\delta)(\kappa-\omega)\bigr)
 =2L\bigl(\omega-\delta(2\kappa-\omega)\bigr)>L
\]
by the fourth hypothesis. So there are primes in $(V,U]$ whose squares do not divide $b$. Let $p$ be the
\emph{largest} one. Every prime $q\in(p,U]$ has $q^2\mid b$, and $q\le U<2V<2p$ because $2\omega<\kappa$.
Let $B$ be the product of these $q^2$. Then $B\mid b$ and $\log B=2(\theta(U)-\theta(p))$.

Consider the integer $b^{dp}Q_p(a/b)=\sum_ic_ia^ib^{dp-i}$. Each of its terms is divisible by $q^{2dp}$ for
each of these $q$. Indeed $q^{2i}$ divides $c_i$ by hypothesis, and $q^2\mid b$ gives $q^{2(dp-i)}\mid b^{dp-i}$.
So $B^{dp}$ divides the integer. Since $p^2\nmid b$, $Q_p(a/b)\ne0$. So $b^{dp}Q_p(a/b)/B^{dp}$ is a nonzero
integer, and
\begin{equation}\label{eq:recycled}
 1\le\Bigl(\frac bB\Bigr)^{dp}\,|Q_p(a/b)| .
\end{equation}
Since $\rho p\le\rho U=\rho\kappa L$, we have $|a/b-\xi|\le e^{-\rho p}$. So $|Q_p(a/b)|\le e^{-\eta p^2}$, and
\eqref{eq:recycled} gives $\eta p^2\le dp\bigl(L-2(\theta(U)-\theta(p))\bigr)$. With
$\theta(U)-\theta(p)\ge U-p-2\delta U$ this is
\[
 \eta p\le d\bigl(L-2U+2p+4\delta U\bigr).
\]
Hence $(\eta-2d)p\le dL-2dU+4d\delta U$. As $\eta<2d$ and $p\le U$, the left side is at least $(\eta-2d)U$. So
$\eta U\le dL+4d\delta U$, that is, $\eta\kappa\le d+4d\delta\kappa$. This contradicts the last hypothesis.
\end{proof}

The largest admissible prime is the right choice. The primes of the window above it have their squares in
$b$, and the divisibility hypothesis recovers these squares from the coefficients. A cruder choice, of any
prime in a short interval that does not divide $b$, loses a factor. The free lower end $V=(\kappa-\omega)L$
allows $\kappa<2$, because the condition $q\le U<2p$ for the recycled primes needs only $U<2V$.

\subsection{The hypotheses for the design $D$}
Take the design $D=(58,9,3,5,3)$ and $d=53$. The polynomials $Q_p$ of the proposition are the $Q_N$ with
$N=p$. Let
\begin{gather*}
 \rho=\Rho,\quad \eta=\EtaTot,\quad \omega=\tfrac35,\quad \delta=2.9\cdot10^{-7},\\
 \kappa=\frac{410}{\rho}=\Kappa\ldots
\end{gather*}
We check the hypotheses.
\begin{itemize}
\item $\deg Q_N\le53N$ by \eqref{eq:QN} and Lemma~\ref{lem:degree}.
\item The nonvanishing. By Proposition~\ref{prop:nonvanishing}, $\Delta_p(a/b)\ne0$ for large $p$ with
$p^2\nmid b$. Hence $Q_p(a/b)\ne0$.
\item The divisibility. By Lemma~\ref{lem:recycling}, $q^{2i}$ divides the $i$-th coefficient of $Q_p$ for all
primes $p<q\le2p$.
\item The smallness. By Theorem~\ref{thm:real}, Proposition~\ref{prop:denominators} and
$\log(53N)!=o(N^2)$, for $|x-G|\le e^{-\rho N}$ and large $N$,
\[
 |Q_N(x)|\le\exp\bigl(-(\TDelta\ldots-\Aden\ldots-3\varepsilon)N^2\bigr)\le e^{-\eta N^2},
\]
since $\eta=\EtaTot<\EtaExact\ldots$ In the formal proof this is the inequality
$\eta<n_0^2(U_0-\frac12\log2)-\mathsf A$ with $\int\kappa$ replaced by its lower bound $\ChiLean$.
\item The numerical conditions. $\rho\kappa=410$ and $0<\eta<2d=106$. Next $\omega=3/5\le1$ and
$2\omega=1.2<\kappa$. Further $2(\omega-\delta(2\kappa-\omega))=1.199998\ldots>1$. Finally
$d+4d\delta\kappa=\PayL\ldots<\eta\kappa=\PayR\ldots$
\end{itemize}
Proposition~\ref{prop:criterion} therefore gives $|G-a/b|>b^{-\rho\kappa}=b^{-410}$ for all integers $a$ and all
sufficiently large $b$. This proves Theorem~\ref{thm:main}. The last condition is equivalent to
$410>x=d\rho/(\eta-4d\delta)=\Bx\ldots$

\section{Remarks on the method}\label{sec:remarks}

\subsection{The three ingredients}
The exponent $d\rho/\eta$ has three ingredients. The degree $d=n-g$ comes from the construction. Of
$\rho=\Rho$, the energy cost of the fixed columns $n_0c_2$ is $208.62$, and the weight and the binomial factor
of one column contribute $\Psi_0=\PsiZero$. Bounding the two separately, by their maxima $52.79$ and
$58\log2=40.20$, would give $\rho=\RhoOld$. The decay rate $\eta$ is a small difference of two large numbers.
The rank deficiencies of Section~\ref{sec:content} contribute $22.05$ of $\eta=\EtaTot$, and the signed bound
of Section~\ref{sec:signed} contributes $\Gain$. Since the exponent is inversely proportional to $\eta$, a further
gain of $0.1$ in $\eta$ lowers it by about $1.1$.

\subsection{The determinant of the row functions}
The energy method bounds $|\det\mathsf L|$ by $6^n$, which is $e^{O(N)}$ and invisible at the scale $N^2$. At
the extremal configuration, numerical evaluation for small $N$ suggests that $\log|\det\mathsf L|$ is about
$-28N^2$, so most of the true size is not captured. The bound \eqref{eq:eleven} loses only $e^{O(N)}$, and
Theorem~\ref{thm:heine} is an identity. The loss is in the log-gas bounds of \S\ref{ssec:loggas}, which use fixed
trial measures, and in the requirement that the combined functional be bounded by its maximum over all
configurations. A sharper treatment of the denominator $Z_s(\psi)$ would raise the gain. These numerical values
are not part of the formal proof.

\subsection{Designs}
All the results of Section~\ref{sec:content} are lower bounds for the nullity, and the proofs give the kernel
vectors explicitly. The closed formulas, the counting inequality and the families derived from it are proved
for every design. So are the other arguments of the paper, except the stored subspaces of
Section~\ref{ssec:packets} and the trials, which belong to the design $D$. A different design changes $d$,
$\rho$ and $\eta$ together. We make no claim that $D$ is optimal. Except for one pair of families, the function
$\kappa$ uses no relation between different families. At a prime where windows overlap the exact nullity may
be larger than each bound. An upper bound for the nullity is not needed. More independence between families,
or larger stored subspaces, would increase $\eta$ without changing anything else.

\section{Formalization}\label{sec:formal}

All results of this paper are formalized in Lean~4 \cite{Lean4} with Mathlib \cite{Mathlib}, on top of OpenAI's
formalization of the irrationality of $G$ \cite{OpenAI2026}. The development is available at
\begin{center}
\FormalizationLink.
\end{center}
Theorem~\ref{thm:main} is the theorem
\begin{center}
{\small\texttt{CatalanMeasure.GenD.Signed.catalan\_irrationalityMeasure\_410}}
\end{center}
in \lean{CatalanMeasure/GenD/Signed/Data/Assembly/CellsAll.lean}. Its statement uses only Mathlib. Catalan's
constant is written as the series $\sum_j(-1)^j/(2j+1)^2$, and the conclusion is $|G-a/b|>b^{-410}$ for all
integers $a$ and all sufficiently large $b$. The modules contain no \lean{sorry} and no \lean{native\_decide}.
The command \lean{\#print axioms} checks that only the three standard axioms \lean{propext},
\lean{Classical.choice} and \lean{Quot.sound} are used. The repository also contains a configuration
\lean{comparator.json} for \lean{lake comparator} \cite{Comparator}. It checks that a solution file proves a
challenge statement with these axioms only. The prime number theorem is taken from the PNT+ project
\cite{PNTPlus}. The parts of OpenAI's formalization that are used were converted to Lean's module system and
ported to a recent Lean release. The repository's \texttt{README} and \texttt{PORTING.md} describe the changes.

\subsection{Module layout}
The design-generic part of the development is in the directory \lean{GenD} of \lean{CatalanMeasure}. A design is
data. The generic modules prove every statement for an arbitrary design, and the instance files contain
numerals only.
\begin{description}[leftmargin=1.5em,style=nextline]\sloppy
\item[Generic layer] \lean{PL} (piecewise linear profiles and their areas), \lean{PrimeGrowth} (the weighted
prime number theorem \eqref{eq:weighted-pnt}), \lean{Budget} (the scalar budgets), \lean{PayW} and
\lean{Recycled} (Proposition~\ref{prop:criterion}), and \lean{Final} (the structure \lean{DesignData} and the
theorem \lean{measure\_of\_design}, which contain no numeral).
\item[Denominators] \lean{OddProfile}, \lean{TwoAdic}, \lean{DenomExp}, \lean{DenomGrowth},
\lean{Denominator} and \lean{Normalizer}, with \lean{Uniform/Arithmetic} and \lean{Uniform/LedgerW} (the
common denominator built from an arbitrary lower bound for the nullity, and the growth of its logarithm).
\item[Rank deficiencies] \lean{RankT}, \lean{RankU}, \lean{RankU6}, \lean{RankS} and their shared layers for
Theorem~\ref{thm:content}; \lean{Uniform/W/} for the counting inequality (\lean{Core}, \lean{Profile}),
Lucas's step (\lean{Lucas/PStep}), the families of Section~\ref{ssec:families} and the stored data of
Section~\ref{ssec:packets} (\lean{Prof/Packet}, \lean{Charts6/Check}); the profile of $D$ with its area
\lean{gain532} and the bound \lean{gain532\_ge} in \lean{Uniform/Bank532}.
\item[Nonvanishing] \lean{Nonvanishing}, with the residue-table certificates modulo $1009$.
\item[Real place, perturbation] \lean{Analytic/} (the energy method with fixed points, the column expansion
and the quadrature nodes); \lean{BinomVandermonde} (Lemma~\ref{lem:binomvand}) and \lean{Keyhole}
(Lemma~\ref{lem:keyhole}) in \lean{GenD/Signed}; \lean{D58/Budget590} (the budgets at the rates $U$ and $U_0$,
and $\rho=$ \lean{rhoEnd}); \lean{Signed/Data/S1} (the $14$ cells of the bound for $\Psi_0$).
\item[Real place, signed bound] \lean{GenD/Signed}: \lean{PickDet} and \lean{RealPlace}
(Lemma~\ref{lem:pick}, Theorem~\ref{thm:C3}, \eqref{eq:eleven}), \lean{LemmaT}, \lean{LemmaTFree} and
\lean{HalfDisc} (Lemma~\ref{lem:T}), \lean{RatioIdentity} and \lean{BernsteinSzego}
(Theorem~\ref{thm:heine}), \lean{TiltedTangent} and \lean{ArcsineDenominator} (Lemmas~\ref{lem:LGup}
and~\ref{lem:LGlow}), \lean{Spike} and its companions (Theorem~\ref{thm:spike}), \lean{ReflectionEnergy}
and \lean{ArcsinePhantom} (Theorem~\ref{thm:reflection}), \lean{MarkedExpansion}, \lean{SigmaCap} and
\lean{EnergyNear} (\S\ref{ssec:fixedrows}), \lean{LemmaK} and \lean{Brackets} (\S\ref{ssec:class}),
\lean{Interp} and \lean{Phase2Contract} (the interpolation of trials), and \lean{Gate} and \lean{Theorem410}
(the assembly).
\item[Trials and cells] \lean{Signed/Data/Trial} (the three trials), \lean{Signed/Data/Consts} (the constants
of the trials, evaluated in exact arithmetic), \lean{Signed/Data/Cells} (the cell engine and the cells of the
three trials) and \lean{Signed/Data/Assembly} (the records of the cell statements and the final theorem).
\end{description}

\subsection{Certificates}
The proof combines conventional formal proofs with numerical certificates checked by the Lean kernel. They
are the inverses of the base matrices modulo $1009$, the trial sequences and covers of
Sections~\ref{sec:real} and~\ref{sec:signed}, the cells of the bound for $\Psi_0$, and the stored subspaces,
kernel vectors and window conditions of Section~\ref{ssec:packets}. All are exact rational or integer
computations, decided by the kernel without compiled code. The interval arithmetic used to generate the trial
data and the covers \cite{Arb} is not trusted. Lean rechecks every inequality. Five devices keep the
certificates small enough for the kernel.
\begin{itemize}
\item \emph{Shared Taylor bounds.} The cells of a cover are grouped in blocks. For each block the derivatives of
the polynomial part of the field are evaluated once, exactly, at the centre of the block, together with one
bound for the remainder. Every cell of the block obtains its expansion by an exact shift.
\item \emph{Split shifts.} In the covers of Section~\ref{sec:signed}, only the low-order coefficients of the
shifted expansion are computed exactly. The others are bounded by one crude table per block.
\item \emph{Accuracy per term.} The logarithms are evaluated by the series
$\log\frac{1+x}{1-x}=2x\sum_ix^{2i}/(2i+1)$ with an explicit geometric tail. The number of terms is chosen for
each logarithmic term separately, from a table of nodes shared by all the cells.
\item \emph{Residue recurrences.} The tables of moments modulo $1009$ are generated by the one-step recurrences
that define them, once proved correct. The kernel checks the recurrence and the products of the matrices with
their claimed inverses, row by row, in natural numbers.
\item \emph{$2$-adic normalisation.} Where a recurrence or a shift produces integers divisible by large powers
of $2$, the tables store the quotients by these powers, with one lemma relating the two. This also avoids
integers that the kernel's hashing of literals treats as equal.
\end{itemize}

\section*{Disclosure of AI use}

The mathematics, the Lean formalization, the certificate generators and the text of this paper were produced by
AI systems at the request of the author, who posed the problems, directed the work and is responsible for this
paper: Claude Opus~5.5, Claude Sonnet~5.5 and Claude Haiku~5.5 (Anthropic), running in Claude Code, and
GPT-6.1 Sol (OpenAI), running in Codex. The underlying irrationality proof for $G$ and its Lean formalization are
OpenAI's \cite{OpenAI2026}.

\subsection*{Acknowledgements}
We thank the authors of Mathlib and of the PNT+ project, and OpenAI for making its formalization available.
